
% . consistently achieving the same minimal loss from 10 random starting trees without getting trapped in local optima, comparable to state-of-the-art software

% With this scheme, we were able to recover the same optimal tree found by RaXML-NG~\cite{kozlov2019} for four gene datasets.

% Using RaXML-NG~\cite{kozlov2019} for branch length optimisation and tree likelihood calculation, we were able to optimise tree topology based on moves on phylo2vec vectors to recover an optimal tree for four gene datasets. as RaXML-NG using a \gls{gtr}~\cite{tavare1986} model~\cite{penn2025}. % Compared to a Newick string, saving/transferring topologies as a Phylo2Vec can be up to x\% more efficient depending on the floating precision, and y\% for trees with branch lengths as a matrix. Whereas this is not a major concern for everyday phylogenetic tasks, this can constitute some significant improvement for large trees ($\approx$ 10,000 taxa).

% Plot: memory efficacy
% subplot 1: phylo2vec vs. newick with ints
% subplot 2: phylo2mat vs. newick with floats
% subplot 3: thousands of tree topologies with species
% subplot 4: thousnads of trees with BLs with species

% Plot: sampling efficacy